\documentclass[11pt,a4paper]{article}

\usepackage[margin=2.4cm]{geometry}
\usepackage{amsmath,amssymb}
\usepackage{booktabs}
\usepackage{siunitx}
\usepackage{tikz}
\usepackage{xcolor}
\usepackage{listings}
\usepackage{enumitem}
\usepackage[hidelinks]{hyperref}
\usepackage{url}
\newcommand{\code}[1]{\texttt{\hyphenchar\font=`\-\relax #1}}

\sisetup{detect-all}
\newcommand{\us}[1]{\SI{#1}{\micro\second}}
\newcommand{\ms}[1]{\SI{#1}{\milli\second}}
\newcommand{\MB}[1]{\SI{#1}{\mega\byte}}
\newcommand{\bitrev}{\operatorname{bitrev}}

\lstset{basicstyle=\ttfamily\small, columns=fullflexible, keepspaces=true,
        frame=single, framesep=4pt, xleftmargin=6pt, xrightmargin=6pt}

\title{Exact negacyclic polynomial multiplication on the GPU:\\
       the optimizations behind \code{fherma-polymul}}
\author{Working notes}
\date{14 October 2026}

\begin{document}
\maketitle

\begin{abstract}
These notes document the optimization work on our entry to the FHERMA challenge
\emph{Polynomial Multiplication (GPU-Accelerated)}, which asks for the exact
product $c = a\cdot b \bmod (X^N+1,\, q)$ of two polynomials with coefficients
of up to a few thousand bits. The solution uses a residue number system (RNS),
a negacyclic number-theoretic transform (NTT) per residue prime, and a Chinese
remainder (CRT) reconstruction. For each optimization we give the problem, how
it was diagnosed, what was changed and what it measured, including the changes
that did not help. At the benchmark point $N = 32768$, $W = 868$ the time per
multiplication on our development machine went from \ms{3.07} to \us{700}; the
code was then generalized to the point $N = 65536$, $W = 1747$ that now decides
the ranking, where it first took \ms{3.09}; kernel fixes brought it to \ms{2.6}, and
reading and writing the caller's memory directly through HMM to about \ms{2.3}
(Point \#2: about \us{610}).
\end{abstract}

\tableofcontents

%======================================================================
\section{The problem and how it is scored}

\paragraph{Specification.} The kernel
\code{polynomial-multiplication/negacyclic@1.0.0} computes, for two polynomials $a, b$ of degree $< N$ with coefficients in $[0,q)$,
\[
  c \;=\; a \cdot b \bmod (X^N + 1,\; q),
\]
exactly, with every output coefficient canonical in $[0, q)$. $N$ is a power of two;
$q$ is the largest prime below $2^W$ with $q \equiv 1 \pmod{2N}$. A coefficient is
stored as $L = \lceil W/32 \rceil$ little-endian 32-bit limbs, so an operand is an
$N \times L$ array of \code{u32}. The specification sets no upper bound on $N$ or $W$
and lists five points: $(N,W) = (64,32)$, $(1024,512)$, $(8192,256)$, $(32768,868)$
and $(65536,1747)$.

\paragraph{Scoring.} The score is the median wall-clock time of one call of
\code{fherma\_run}, measured from the moment $a$ and $b$ are available in host
memory to the moment $c$ is complete in host memory. \emph{Copies to and from the
GPU are therefore inside the timed window.} Setup (\code{fherma\_init}) is not
timed but may depend only on $N$, $W$, $L$ and $q$. Results must be exact; FHERMA
re-tests on 20 random seeds per point and on non-public inputs.

\paragraph{Benchmark points.} Since 28 September 2026 the ranking is decided at
\textbf{Point \#1}, $N = 65536$, $W = 1747$, $L = 55$; the original point,
$N = 32768$, $W = 868$, $L = 28$, is now \textbf{Point \#2} and is shown beside it.
An entry with no Point \#1 result ranks behind every entry that has one.

\paragraph{Hardware.} All measurements in these notes come from our development
machine, a shared QEMU virtual machine with 96 virtual CPUs (x86-64-v2, no AVX) and
two NVIDIA L40S cards (\code{sm\_89}, \SI{46}{\giga\byte}, PCIe Gen4 x16). The
ranked runs use an NVIDIA RTX PRO 6000 (\code{fherma-gpu-rtx6000}), a different
architecture with a faster PCIe link, so absolute times are not comparable; the
relative effects of host-side changes in particular may differ there.

\paragraph{Data volume.} At Point \#2 an operand is $32768 \times 28 \times 4$ bytes
$=$ \MB{3.67}, so each call moves \MB{11} across PCIe (two operands in, one out);
at Point \#1 an operand is \MB{14.4} and a call moves \MB{43}. On our Gen4 link,
pinned DMA runs at about \SI{24.5}{\giga\byte\per\second}, which puts a floor of
roughly \us{450} under Point \#2 and \ms{1.75} under Point \#1 before any compute.

%======================================================================
\section{Algorithm}

\subsection{Residue number system}
Choose $K$ primes $p_j < 2^{31}$ with $p_j \equiv 1 \pmod{2N}$ (so that each has a
primitive $2N$-th root of unity $\psi_j$), and let $P = \prod_j p_j$. The integer
(not reduced) negacyclic product $\tilde c = \sum_{i+k=n} a_i b_k - \sum_{i+k=n+N} a_i b_k$
satisfies $|\tilde c_n| < N q^2$. If $P > 2^{20} \cdot 2N q^2$, then $\tilde c$ is
determined by its residues modulo the $p_j$, with a margin of 20 bits that makes the
floating-point rounding below exact. The number of primes needed is
\[
  K \approx \frac{2W + \log_2 N + 1 + 20}{31},
\]
which gives $K = 58$ at Point \#2 and $K = 115$ at Point \#1. The primes are taken
from just below $2^{31}$ downwards in steps of $\max(2N, 2^{16})$; all of them lie
above $2^{30.8}$.

\subsection{Pipeline}
\begin{enumerate}[leftmargin=*]
\item \textbf{Decompose.} Each coefficient's $L$ limbs become its residue modulo each
  prime: $x \bmod p = \sum_{l} x_l \cdot (2^{32l} \bmod p) \bmod p$, with the
  constants $2^{32l} \bmod p$ precomputed (with Shoup companions,
  \S\ref{sec:shoup}). Output: two $K \times N$ arrays of residues.
\item \textbf{Forward NTT} per prime, Cooley--Tukey, with twiddles
  $\psi^{\bitrev(i)}$ so that the negacyclic twist is folded into the transform and
  no pre- or post-multiplication by powers of $\psi$ is needed.
\item \textbf{Pointwise product} in the NTT domain, multiplied by the constant
  $N^{-1} \cdot (P/p_j)^{-1}$, so that the inverse transform directly yields
  $y_j = \tilde c \cdot (P/p_j)^{-1} \bmod p_j$.
\item \textbf{Inverse NTT} per prime, Gentleman--Sande, natural order out.
\item \textbf{CRT.} With $v = \operatorname{round}\big(\sum_j y_j / p_j\big)$,
  computed in double precision, the centred value is
  $\tilde c = \sum_j y_j (P/p_j) - vP$, and modulo $q$
  \[
    c \;=\; \sum_j y_j \cdot M_j \;+\; (-vP \bmod q) \pmod q,
    \qquad M_j = (P/p_j) \bmod q,
  \]
  where both $M_j$ and the $K+1$ possible values of $-vP \bmod q$ are precomputed.
  The final reduction modulo $q$ uses a quotient estimated in double precision from
  the top limbs, followed by at most two subtractions.
\end{enumerate}

\subsection{Modular arithmetic}\label{sec:shoup}
All residue arithmetic is on 32-bit words with $p < 2^{31}$, so a lazily reduced
value below $2p$ still fits in a word.
\begin{itemize}[leftmargin=*]
\item \textbf{Shoup multiplication} by a constant $w$: with the precomputed
  companion $w' = \lfloor w \cdot 2^{32} / p \rfloor$,
  $r = a w - \lfloor a w' / 2^{32} \rfloor \cdot p \pmod{2^{32}}$ lies in $[0, 2p)$
  for any 32-bit $a$; one conditional subtraction finishes it. Three integer
  multiplies, no 64-bit product.
\item \textbf{Montgomery multiplication} $\operatorname{mont}(a,b) = a b \cdot 2^{-32} \bmod p$
  for the pointwise product of two variables.
\end{itemize}

\subsection{GPU schedule}
\begin{itemize}[leftmargin=*]
\item \textbf{Four-step NTT split.} The $\log_2 N$ stages are split at row length
  $S = 2^{\lceil \log_2 N / 2 \rceil}$: stages whose butterflies span a distance
  $\ge S$ pair elements of the same column (stride $S$), the others pair elements
  within a contiguous row. Each pass keeps a tile of $T = 2048$ residues in shared
  memory for all its stages. At $N = 32768$: $S = 256$, $R = N/S = 128$, and a
  column-pass block holds $C = T/R = 16$ columns.
\item \textbf{Kernels per call.} \code{k\_decompose} (per input chunk),
  \code{k\_ntt\_cols} and \code{k\_ntt\_rows\_fwd} on $a$, \code{k\_ntt\_cols} on $b$,
  \code{k\_ntt\_rows\_mul} (forward rows of $b$, pointwise product, inverse rows),
  \code{k\_ntt\_cols} (inverse), and \code{k\_crt} (per output chunk).
\item \textbf{Streams.} Host-to-device copies run on one stream, kernels on a
  second, device-to-host copies on a third, chained by events, so that copying
  chunk $i+1$ overlaps decomposing chunk $i$, and copying output chunk $i$ overlaps
  the CRT of chunk $i+1$.
\item \textbf{Shared schedule.} Tile geometry, index maps and per-phase steps live in
  one header used by both the GPU kernels and a CPU emulation, which runs the same
  schedule block by block and phase by phase and checks that no NTT stage touches
  an element twice. The CPU build therefore tests the GPU's indexing, not just its
  arithmetic.
\end{itemize}

%======================================================================
\section{How we measured}

\paragraph{Correctness.} Every change was checked with \code{run\_tests.sh}, which
generates test points with Python oracles (an exact $O(N^2)$ reference for small $N$,
a GMP Kronecker-substitution reference for large $N$) and verifies the outputs
bit for bit. The matrix now has 21 points: every NTT geometry, every limb edge
($L = 1$, $q$ filling its limbs, $q$ one bit into a new limb, $L = 64$), the five
points of the specification, $N = 2^{17}$, and an all-$(q-1)$ worst case in
case 0 of each large point. Both the CPU build and the GPU build run it. NVIDIA's
\code{compute-sanitizer} (memcheck, racecheck, synccheck) was run after every
kernel change.

\paragraph{Fuzzing.} \code{fuzz.py} adds 61 more points (every 32-bit width boundary up
to 2048, $N$ from 2 to $2^{22}$, 30 random sizes) with six adversarial input patterns
each (all $q-1$, random, zero, the $x^{N-1}\cdot x^{N-1}$ wrap, the identity, alternating),
under every copy mode and twiddle path and on the CPU build, against a GMP reference
that is itself checked against the schoolbook product for small $N$: no failures.
Out-of-range sizes are refused with a message.

\paragraph{Performance.} The development server is shared, and other users' jobs
moved wall-clock medians by \us{50} or more, with outliers up to \ms{1.3}. Three
practices made the measurements usable:
\begin{enumerate}[leftmargin=*]
\item \textbf{Interleaved A/B runs}: alternate the builds or settings being compared
  over 5--20 rounds of 21 cases, and compare medians and lower quartiles of the
  per-round medians.
\item \textbf{Kernel times and the GPU span from Nsight Systems}: the duration of each
  kernel and the span from a call's first transfer to its last are almost
  unaffected by host noise.
\item \textbf{Elimination}: build scratch variants that skip one part of a kernel
  (the stores, the arithmetic, the table loads) and attribute the time by difference.
  Combined with instruction counts from the SASS (\code{cuobjdump}), this found
  the real bottleneck every time our first guess was wrong.
\end{enumerate}

%======================================================================
\section{Optimizations}

Unless stated otherwise, numbers are for Point \#2. Each subsection gives the
problem, the diagnosis, the change and the measured effect.

\subsection{CRT: from one thread per coefficient to prime groups}

\paragraph{Problem.} The CRT kernel took \us{47} in one launch over all 32768
coefficients, but about \us{38} per launch \emph{regardless of launch size} when
the output was split into chunks, so the default 8 chunks cost \us{308}. One
coefficient was handled by one thread, in blocks of 64 threads: a chunk of 4096
coefficients was 64 blocks of two warps on a 142-SM GPU.

\paragraph{Diagnosis.} Removing each part of the kernel in turn gave, per chunk:
\begin{center}
\begin{tabular}{lr}
\toprule
part of \code{k\_crt} & cost per chunk \\
\midrule
products ($58 \times 28$ multiply-accumulates) & \us{13.5} \\
loading $Y$, storing $c$, launch & \us{12} \\
rounding estimate $\sum_j y_j/p_j$ & \us{2.3} \\
final reduction modulo $q$ & \us{2} \\
\bottomrule
\end{tabular}
\end{center}
No single part was slow: every phase was waiting on latency, because there was far
too little parallel work per launch.

\paragraph{Change.} (1) Each block first stages its tile of $Y$ ($K$ residues for
each of its coefficients) into shared memory in one coalesced pass, with all loads
in flight together. (2) The $K$ primes are split into 4 groups, and each coefficient
gets 4 threads, each computing a partial sum over its group, carried down to $L+2$
words. A last phase adds the 4 partials and $-vP \bmod q$ and reduces modulo $q$.
Blocks are 32 coefficients $\times$ 4 groups $= 128$ threads; a warp is one prime group
and 32 consecutive coefficients, so its loads of the table $M$ are uniform.

\paragraph{Result.} \us{308} $\to$ \us{92} at 8 chunks; \us{47} $\to$ \us{38} in one
launch. The staging step alone (without the groups) gave only
\us{308} $\to$ \us{232}: latency of global loads was part of the problem, the lack of
parallelism the larger part. End to end this was worth only 1--2\,\% at the time,
because the host side dominated (next section).

\subsection{Host side: a copy pool instead of page registration}

\paragraph{Problem.} With the GPU work at about \ms{1.1} of a \ms{2.1} call,
roughly half the call was spent on the host. The inputs arrive in freshly
allocated \code{std::vector}s, and the output must be returned as a new vector.

\paragraph{Diagnosis.} A host-only benchmark measured each step for one
\MB{3.67} operand on this VM:
\begin{center}
\begin{tabular}{lr}
\toprule
step & median \\
\midrule
\code{cudaHostRegister} + unregister of the caller's buffer & \us{400} \\
\code{memcpy} into a pinned buffer, 1 thread & \us{300} \\
\code{memcpy} into a pinned buffer, 16 threads (cache-hot) & \us{13} \\
zero-filling a new output vector (\code{resize}) & \us{237} \\
pinned DMA, either direction & \us{150} \\
pageable \code{cudaMemcpy}, either direction & \us{357} \\
\bottomrule
\end{tabular}
\end{center}
Registering the caller's pages, the fastest mode before, was the most expensive
way to reach the GPU; one-thread copies through pinned memory were next.

\paragraph{Change.} A pool of \code{std::thread} workers copies each input chunk into
a pinned staging buffer, then the chunk's DMA and decompose are enqueued. The copy
is cut into \SI{64}{\kilo\byte} pieces that the caller and the workers claim from a
shared counter. The counter packs a job number with the piece index and pieces are
claimed by compare-and-swap, so a late worker can never claim a piece of a newer
job with an older job's parameters, and the caller never waits for a worker that is
not awake. The output is copied out of a pinned buffer by the same pool, chunk by
chunk, as each chunk lands. A first version used static slices per worker; it had a
shutdown hang (spinning workers never checked the quit flag), found by an
experiment and fixed.

\paragraph{Result.} \us{2104} $\to$ about \us{880} per call (best setting before
$\to$ new default), a factor of 2.4.

\subsection{Chunking: fewer, shaped chunks}

\paragraph{Observation.} After the copy pool, more chunks made things \emph{worse}:
8 input chunks cost about \us{250} more than 2. We first attributed this to about
\us{25} of CUDA API calls per chunk. That explanation was wrong: with zero-copy
output (\S\ref{sec:zerocopy}), where a chunk is only one kernel launch and one event,
extra chunks still hurt. The cost is a fixed overhead per parallel copy in the pool.

\paragraph{Changes and results.}
\begin{itemize}[leftmargin=*]
\item Default pipeline depth reduced from 8 chunks to 2.
\item Input and output chunking were decoupled, and the chunks of each operand can
  be given relative weights (settings in Appendix~\ref{app:settings}). An output split of $1{:}2{:}1$ lets the first output
  DMA start after a smaller first CRT chunk and leaves less to copy after the last
  DMA. Over 10 interleaved rounds: \us{715} $\to$ \us{700} median, lower quartile
  \us{710} $\to$ \us{695}. Now the default.
\item Uneven \emph{input} chunks (a small first chunk of $a$, a small last chunk of
  $b$) did not help; a $3{:}1$ or $7{:}1$ split of $b$ made the GPU span worse
  (\us{625} and \us{647} against \us{591}--\us{610}), because the larger first chunk
  delays $b$'s first transfer.
\end{itemize}

\subsection{Zero-copy decompose and CRT}\label{sec:zerocopy}

\paragraph{Idea.} Map the pinned staging buffers into the device address space so
that \code{k\_decompose} reads its limbs directly over PCIe and \code{k\_crt} writes
its limbs directly to host memory, removing the separate DMA, event and stream
hop of every chunk.

\paragraph{Result.} Over 5 interleaved rounds, mapped input saved about
\us{15}--\us{20} (2\,\%), well inside the \us{790}--\us{940} spread between rounds.
Kept as an opt-in mode (\code{FHERMA\_IN/OUT=mapped}); the default stays staged.
It may behave differently on the reference machine.

\subsection{Worker wake-up: spinning between calls}

\paragraph{Problem.} The copy workers sleep on a condition variable between calls.
Timing each input copy showed the first copy of a call taking \us{240}--\us{270}
instead of \us{60}--\us{110}: waking sleeping threads in this VM costs about
\us{150}--\us{180}. The harness leaves about \ms{13} between calls (it writes the
previous answer and reads the next inputs from files).

\paragraph{Change.} \code{FHERMA\_SPIN\_MS} (default \SI{20}{\milli\second}): after a
call, the workers keep spinning for that long before sleeping. A test without the
\code{pause} instruction in the spin loop (to rule out the hypervisor descheduling
spinning virtual CPUs) made no difference.

\paragraph{Result.} Over 5 interleaved rounds: \us{860} $\to$ \us{760} (staged) and
\us{840} $\to$ \us{763} (mapped); every spinning round beat every sleeping round. The
cost is 15 threads busy-waiting through the gap between calls; whether the grader
allows that is an open question.

\subsection{Zero-filling the output on a helper thread}

\paragraph{Problem.} The output vector must be sized (\code{resize}, a zero fill
of about \us{350} here) before chunks can be copied into it. That happened after all
GPU work was enqueued, about \us{300} into the call, so it ended near \us{660},
after the first output chunk had landed (about \us{620}): it was on the critical path.

\paragraph{Change.} A dedicated helper thread (spinning and sleeping like the pool)
zero-fills the output from the start of the call, while the calling thread copies
the inputs. A guard object waits for it on every exit path, since it writes into
the vector being returned.

\paragraph{Result.} \us{736} $\to$ \us{708} over 6 interleaved rounds (5 consistent,
one outlier).

\subsection{Overlapping $a$'s forward NTT with $b$'s transfer}

\paragraph{Problem.} The forward transforms of $a$ and $b$ ran together after $b$ had
fully arrived, but $a$'s does not need $b$.

\paragraph{Change.} The forward transform of $a$ (column pass and a new row-only
kernel \code{k\_ntt\_rows\_fwd}) is enqueued right after $a$'s decompose, so it runs
while $b$ is still crossing PCIe; \code{k\_ntt\_rows\_mul} gained a mode that skips
$a$'s forward stages. The CPU emulation follows the same order.

\paragraph{Result.} GPU span \us{622} $\to$ \us{591} (median of 21 calls; the minimum
improved by the same amount). The column pass on $b$ alone also costs less than the
old pass on both operands (\us{16} against \us{22}).

\subsection{Decompose: register blocking}

\paragraph{Problem.} \code{k\_decompose} took \us{20.6} for half an operand, and the
last call of it is on the critical path.

\paragraph{Diagnosis.} Two guesses failed first: making $L$ a compile-time constant
(unrolled loops; \us{19.8}) and lazy reduction (\us{19.4}), so neither latency nor
arithmetic instruction count was the limit. Elimination showed the arithmetic phase
was essentially the whole cost (\us{19.4} full, \us{2.8} without the arithmetic,
\us{19.1} without the stores). The SASS of \code{k\_decompose<28>} showed, per residue,
58 global loads, 28 shared loads, 101 \code{IMAD} and 37 \code{IADD3}: \textbf{two
table loads per term}, the two words $2^{32l} \bmod p$ and its Shoup companion,
reloaded by every thread although a whole warp shares one prime.

\paragraph{Change.}
\begin{itemize}[leftmargin=*]
\item \textbf{Register blocking}: each thread computes 4 coefficients for the same
  prime (lanes take coefficients $\ell, \ell+32, \ell+64, \ell+96$), so each table
  entry loaded serves 4 terms, and every store is still 32 consecutive coefficients.
\item \textbf{Lazy reduction}: each term is reduced only to $[0, 2p)$, terms are summed
  in 64 bits (below $2^{37}$ for $L \le 32$), and the sum $h \cdot 2^{32} + \ell$ is
  reduced once as $\ell \cdot 1 + h \cdot (2^{32} \bmod p)$, using entries 0 and 1 of
  the existing table.
\item \textbf{Block size} 512 threads instead of 256, for enough warps per SM.
\end{itemize}

\paragraph{Result.}
\begin{center}
\begin{tabular}{lr}
\toprule
version & time (half operand) \\
\midrule
original & \us{20.6} \\
$L$ templated & \us{19.8} \\
\quad + lazy reduction & \us{19.4} \\
\quad + 4 coefficients per thread, 256 threads & \us{15.4} \\
\quad + 2 coefficients per thread, 512 threads & \us{13.8} \\
\quad + 4 coefficients per thread, 512 threads & \textbf{\us{12.2}} \\
\bottomrule
\end{tabular}
\end{center}

\subsection{NTT twiddles: one load instead of two}

\paragraph{Diagnosis.} By elimination, \code{k\_ntt\_rows\_mul} (\us{31.8}) split into
about \us{14} of data movement, \us{8.4} of butterfly arithmetic and \us{9.4} of
twiddle loads; \code{k\_ntt\_cols} (\us{15.9}) into \us{10.1}, \us{2.5} and \us{3.3}.

\paragraph{Change.} Each twiddle and its Shoup companion are stored side by side
(\code{struct alignas(8) Twiddle}), so a butterfly fetches both with one 8-byte load.

\paragraph{Result.} Only about \us{1} per kernel (\us{31.8} $\to$ \us{30.8}, \us{15.9}
$\to$ \us{15.3}). The twiddle cost is bytes, not instructions: a row-pass block needs
about 2048 distinct twiddles per transform, so \code{k\_ntt\_rows\_mul} reads about
\MB{30} of twiddles per call from L2. Kept because it is simpler and not slower.

%======================================================================
\section{Generalizing to Point \#1}

The original code accepted only $N \le 2^{15}$, $L \le 32$ and $K \le 80$, and stopped
in \code{fherma\_init} at Point \#1. Changes:
\begin{itemize}[leftmargin=*]
\item Primes $p \equiv 1 \pmod{\max(2N, 2^{16})}$: the same set as before for
  $N \le 2^{15}$, a valid set for larger $N$.
\item $N$ up to $2^{22}$ (beyond that, a row would not fit the 2048-element tile),
  $L$ up to 64 ($W \le 2048$), $K$ up to 160; anything else is rejected with an
  explicit message.
\item The CRT block's shared memory grows with $K$ and $L$ (about \SI{52}{\kilo\byte}
  at Point \#1) and is now allocated dynamically, with the opt-in attribute set past
  \SI{48}{\kilo\byte} and a clear error if the device cannot provide it.
\end{itemize}
All 21 test points pass on the CPU and the GPU, including Point \#1 with all-$(q-1)$
inputs, $W = 2048$ against the exact oracle and $N = 2^{17}$; the sanitizers are clean
at $L = 55$ and $L = 64$. Register allocation shows no spills at $L = 55$.

The submission was also built and run inside FHERMA's challenge image
\code{fherma/cupqc:0.6.0} (the digest shown on the leaderboard) from a clean copy of
only the submitted files, with the platform's exact build command: no warnings, all
points verified, and the same speed as the native build.

%======================================================================
\section{Optimizations for Point \#1}

At Point \#1 the first working version spent \us{707} in kernels between the last
input transfer and the first output transfer (against \us{95} at Point \#2), with two
kernels scaling far worse than the data, which grows by a factor of about 4 ($K$ from 58
to 115, $N$ doubled):
\begin{itemize}[leftmargin=*]
\item \code{k\_ntt\_rows\_mul}: \us{31} $\to$ \us{335}. The full twiddle tables are
  $16 K N$ bytes $=$ \MB{121}, more than the L40S's \SI{96}{\mega\byte} L2 (and, with the
  data, than the RTX PRO 6000's \SI{128}{\mega\byte}), so they stream from DRAM.
\item CRT: \us{12} $\to$ about \us{170} per chunk. Its shared memory had grown to
  \SI{52}{\kilo\byte} per block, which allows one 4-warp block per SM.
\end{itemize}

\subsection{Twiddles generated from two small tables}

\paragraph{Idea.} A twiddle with index $i$ is $\psi^{e}$ with $e = \bitrev(i)$.
Splitting $e = e_h 2^s + e_\ell$ with $s = \lceil \log_2 N / 2 \rceil$,
\[
  \psi^{e} = \psi^{e_h 2^s} \cdot \psi^{e_\ell},
\]
so two per-prime tables of $2^{\log_2 N - s}$ and $2^s$ entries suffice: 256 + 256 at
Point \#1, about \SI{0.5}{\mega\byte} for all primes and both directions, against
\MB{121}. Stored in Montgomery form ($x \cdot 2^{32} \bmod p$), the twiddle is one
Montgomery product of two table entries,
$\operatorname{mont}(\psi^{e_h 2^s} 2^{32}, \psi^{e_\ell} 2^{32}) = \psi^e 2^{32}$, and
applying it is a second one, $\operatorname{mont}(y, \psi^e 2^{32}) = y \psi^e$; no Shoup
companion is needed.

\paragraph{Implementation.} The butterflies take a \emph{twiddle source} with one
method, \code{mul(y, i)} $= y \cdot w_i \bmod p$: \code{TableTwiddles} reads the full
Shoup table, \code{GenTwiddles} builds $w_i$ as above. The row kernels are templated on
the source. The column pass only ever uses indices below $2^{\lfloor \log_2 N/2 \rfloor}$,
so it always reads a small head of the full table instead (\SI{235}{\kilo\byte} at
Point \#1). The full tables are built only when the table path is chosen, which also
saves \MB{121} of device memory and part of the setup time at Point \#1.

\paragraph{Result.} Kernel times (Nsight Systems):
\begin{center}
\begin{tabular}{lrrrr}
\toprule
 & \multicolumn{2}{c}{Point \#1} & \multicolumn{2}{c}{Point \#2} \\
kernel & table & generated & table & generated \\
\midrule
\code{k\_ntt\_rows\_mul} & 336 & \textbf{132} & \textbf{31} & 38 \\
\code{k\_ntt\_rows\_fwd} & 120 & \textbf{64} & \textbf{17} & 20 \\
\code{k\_ntt\_cols} (head table) & 65 & \textbf{55} & 15 & 15 \\
\bottomrule
\end{tabular}
\end{center}
(times in \si{\micro\second}). Generated twiddles win when the tables fall out of L2 and
lose when they fit, so the choice is automatic: generated when $16KN$ exceeds
\MB{64} (Point \#1), table otherwise (Point \#2), overridable with
\code{FHERMA\_TWIDDLES}. The test matrix passes on both paths. GPU span at Point \#1:
\ms{2.58} $\to$ \ms{2.32}.

\subsection{CRT split over limb ranges}

\paragraph{Change.} Instead of splitting each coefficient's sum over groups of primes,
each of the block's 4 warps owns a range of output limbs: warp $g$ accumulates the 64-bit
columns of limbs $[g\cdot\textit{span},\, (g+1)\cdot\textit{span})$, with
$\textit{span} = \lceil L/4 \rceil$, of $\sum_j y_j M_j$ over \emph{all} primes, plus the
rounding sum $\sum y_j/p_j$ over the primes $j \equiv g \pmod 4$. Lane $i$ takes
coefficient $i$ of the block, so its $y_j$ loads are 32 consecutive coefficients of one
prime (coalesced, read straight from global memory) and its loads of $M$ are uniform
across the warp. Each partial is carried down to $\textit{span}+2$ words; a combine step
per coefficient adds the 4 partials and $-vP \bmod q$, reduces modulo $q$, and writes the
result over the same shared-memory row (odd row stride, so no bank conflicts). Shared
memory no longer depends on $K$: about \SI{10}{\kilo\byte} at $L = 55$ instead of
\SI{52}{\kilo\byte}.

\paragraph{Tuning.} Kernel time per chunk:
\begin{center}
\begin{tabular}{lrr}
\toprule
variant & Point \#2 & Point \#1 \\
\midrule
old: prime groups, $Y$ staged in shared memory & 11.9 & about 170 \\
limb ranges, 2 coefficients per thread & 31.5 & 53.6 \\
limb ranges, 1 coefficient per thread & 15.9 & 50.0 \\
\quad + prime loop unrolled by 4 & \textbf{13.1} & \textbf{50.1} \\
\quad + unrolled by 8 & 18.3 & 50.3 \\
\bottomrule
\end{tabular}
\end{center}
(times in \si{\micro\second}). Two coefficients per thread halve the number of blocks,
and each iteration of the prime loop starts with a global load that has little
arithmetic to hide behind; unrolling by 4 keeps several of those loads in flight. At
Point \#2 only the small first chunk is on the critical path, so the \us{1} lost there is
negligible.

\subsection{Host overhead at Point \#1}
The wall clock exceeded the GPU span by about \us{580}. The host timeline shows the input
copies into pinned memory saturating at about \SI{29}{\giga\byte\per\second} (\ms{1} for
\MB{28.8}) from 16 threads on, about the PCIe rate, so the first transfer waits for the
whole first chunk (about \us{265}), and the last output chunk is copied out after the last
transfer. Nothing tried helped measurably:
\begin{itemize}[leftmargin=*]
\item more copy threads (24, 32, 48): copy time unchanged;
\item seven chunk shapes, including a small first chunk of $a$ and small last output
  chunks: all within \ms{2.60}--\ms{2.81}, inside the noise;
\item starting the output zero-fill only after the first input chunk, to free memory
  bandwidth for it: worse (\ms{2.69} $\to$ \ms{3.08}), the fill then lands on the critical path.
\end{itemize}
On this VM the host side is bound by memory bandwidth; it may well be less so on the
reference machine.

\subsection{A container-only slowdown}
Built and run inside FHERMA's image, Point \#1 is 0.3--\ms{1} slower than natively (Point
\#2 only about \us{50}). The host timeline inside the container shows the helper's zero-fill
of the \MB{14.4} output vector taking about \ms{3} instead of \ms{1.3}. Page faults are the
same in both. The image's glibc 2.39 zero-fills blocks above \SI{2}{\kilo\byte} with
\code{rep stosb}, while the host's glibc 2.43 uses non-temporal stores above
\SI{256}{\kilo\byte}; this VM's CPU model advertises neither ERMS nor FSRM, the features
that make \code{rep stosb} fast. The slowdown is therefore probably specific to our VM;
a fix that does not depend on the C library would be to zero-fill the next call's output
vector in the background between calls, which moves work out of the timed window and is
not done.

%======================================================================
\section{Getting closer to the transfer floor}

Every output coefficient depends on every input coefficient, so the output cannot start
before the whole input has arrived: the floor is the input transfer plus the output
transfer. At Point \#1 the wall clock (\ms{2.6}) exceeded that floor (\ms{1.8} on our
Gen4 link) by about \us{400}--\us{500} of host copies at the edges of the call and
\us{370} of kernels between the last input byte and the first output byte.

\subsection{Reading and writing the caller's memory through HMM}
The driver here runs in HMM mode (heterogeneous memory management: GPU kernels can access
ordinary pageable host memory). With the advice ``preferred location: host, accessed by:
GPU'' (\code{cudaMemAdvise}), a kernel reads or writes a pageable buffer over PCIe without
migrating its pages:
\begin{center}
\begin{tabular}{lrr}
\toprule
14.4 MB (one Point \#1 operand) & time & rate \\
\midrule
pinned DMA (reference) & \us{593} & \SI{24.3}{\giga\byte\per\second} \\
kernel reads the pageable vector, no advice (pages migrate) & \ms{32} & \SI{0.4}{\giga\byte\per\second} \\
kernel reads it, advised to stay in host memory & \us{684} & \SI{21.1}{\giga\byte\per\second} \\
prefetch (migrate) to the GPU, then read & \ms{6.3} & \SI{2.3}{\giga\byte\per\second} \\
kernel writes a pageable vector, advised to stay in host memory & \us{732} & \SI{19.7}{\giga\byte\per\second} \\
previous output path (DMA, then one-thread copy) & \ms{2.8} & \SI{5.1}{\giga\byte\per\second} \\
\bottomrule
\end{tabular}
\end{center}
So in the new default mode (\code{hmm}) the decompose kernel reads the limbs straight from
the caller's input vectors, and the CRT writes the result straight into the output vector
(after the helper's zero-fill): there are no host staging copies at all. The output stays
in host memory, so nothing is deferred past the end of the call. Where the device cannot
access pageable memory the code falls back to the staged path.

Two consequences followed. First, in this mode $b$'s decompose kernels \emph{are} $b$'s
transfer, so $a$'s forward NTT moved to its own stream, with lower priority than the
decompose stream (otherwise $b$'s first decompose started \us{46} late, behind the NTT's
blocks). Second, the decompose kernel now loads 16 bytes at a time for full blocks, which
raised its read rate from \SI{18.4}{} to about \SI{20}{\giga\byte\per\second}.

\subsection{Chunking in HMM mode}
With no host copy to overlap, every chunk boundary only leaves the end of a kernel
underusing the GPU, and here the kernels are the transfers. GPU span at Point \#1:
\begin{center}
\begin{tabular}{lr}
\toprule
input chunks / output chunks & GPU span \\
\midrule
2 / $1{:}2{:}1$ (staged defaults) & \us{2404} \\
8 / $1{:}2{:}1$ & \us{2705} \\
1 / $1{:}2{:}1$ & \us{2323} \\
2 / 1 & \us{2304} \\
\textbf{1 / 1} & \textbf{\us{2242}} \\
\bottomrule
\end{tabular}
\end{center}
One chunk each way is the default in \code{hmm} mode.

\paragraph{Result.} Interleaved medians, native: Point \#1 \us{2616} $\to$ about \us{2300};
Point \#2 \us{684} $\to$ about \us{610}. Host overhead fell from about \us{400}--\us{500} to
about \us{60} at Point \#1. All 21 test points pass in every copy mode; sanitizers clean.

\subsection{Tried and not adopted}
\begin{itemize}[leftmargin=*]
\item \textbf{A pinned memory pool behind a replaced global \code{operator new}}, so that the
  harness's own input and output vectors would be pinned and DMA-able at the full
  \SI{24.5}{\giga\byte\per\second}. Built only as a scratch experiment: it deadlocked inside the
  CUDA runtime, even with a re-entrancy guard and with the pool enabled only after
  initialization. It would also change how the harness allocates its buffers, so it would need
  FHERMA's approval. Expected gain over \code{hmm}: the DMA-versus-HMM rate difference, about
  \us{250}--\us{300} at Point \#1.
\item \textbf{A single-pass NTT per prime using thread-block clusters} (distributed shared memory),
  which would cut the memory passes after the last input byte from three to one. Clusters need
  \code{sm\_90} or newer; the L40S cannot run them, so this was not attempted.
\item \textbf{A measurement on FHERMA's hardware}: it requires pushing the code to a repository and
  registering the commit from the participant's account.
\end{itemize}

\subsection{The zero-fill inside the challenge image}
Inside \code{fherma/cupqc:0.6.0} on this VM, every mode waits on the output vector's
zero-fill, which takes 1.5--\ms{2.9} there depending on load, against about \ms{1.3}
natively (the image's glibc 2.39 \code{memset} against the host's 2.43). In \code{hmm}
output mode the CRT has to wait for it, so measured in the image (three rounds) the modes
overlap: staged \ms{3.1}--\ms{3.7}, \code{hmm} 2.5--\ms{3.6}. This is probably specific to
our VM's old CPU model, but it is unverified for FHERMA's machine. The only complete fix
found, preparing the next call's zeroed output vector between calls, moves work out of the
timed window and is not done.

\subsection{Further attempts}
\begin{itemize}[leftmargin=*]
\item \textbf{Huge pages for the HMM buffers.} With the advice calls outside the timing (as
  the code caches them), a streaming kernel already reaches \SI{23.9}{\giga\byte\per\second}
  reading and \SI{25}{\giga\byte\per\second} writing pageable memory with \SI{4}{\kilo\byte}
  pages; collapsing the range to \SI{2}{\mega\byte} pages changed nothing. The shortfall
  against DMA is in our kernels, which stop reading or writing while they compute.
\item \textbf{A separate streaming copy kernel} (\code{stream} mode) moves data between the
  caller's memory and device memory at full speed, but it starves the decompose of SMs
  (\us{902} instead of \us{72} for the first chunk): slower overall (GPU span \us{2452}
  against \us{2245}). Kept as an option.
\item \textbf{Two decompose streams.} Alternating decompose chunks between two equal-priority
  streams lets each chunk start while the previous one finishes: two chunks now cost the same
  as one, eight cost \us{2372} instead of \us{2705}.
\item \textbf{Column-group streaming} of $b$'s column pass and of the inverse column pass
  with the CRT was not built: it needs about eight groups per operand, and eight decompose
  chunks still cost about \us{120} more than one, against about \us{50} of hidden work.
\item \textbf{A software-pipelined decompose} (persistent blocks, two shared-memory tiles,
  \code{cp.async} loads for the next tile while the current one is computed): correct, but
  slower (\us{895} per Point \#1 operand against \us{667}): the two tiles take
  \SI{56}{\kilo\byte}, so one block per SM, and word-by-word asynchronous copies cost more
  instructions than they save.
\item \textbf{Choosing the copy mode at setup.} Which mode wins depends on the machine, so
  \code{fherma\_init} times the candidates on dummy operands (seven runs each, evicting the
  caches between runs to mimic the harness's file I/O) and keeps the default unless another
  mode is at least 5\,\% faster.
\end{itemize}

%======================================================================
\section{Where the time goes now}

\subsection{Point \#2}
Wall clock about \us{700} median per call; GPU span about \us{590}. One call:
\begin{center}
\begin{tabular}{rp{0.74\textwidth}}
\toprule
time (\si{\micro\second}) & activity \\
\midrule
0--326 & four \MB{1.84} host-to-device transfers, back to back; $a$'s decompose and
          forward NTT hidden under them \\
326--434 & about \us{95} of kernels: last decompose, column pass on $b$,
          \code{k\_ntt\_rows\_mul}, inverse column pass, first CRT chunk \\
434--582 & device-to-host transfers of \MB{0.92}, \MB{1.84} and \MB{0.92}, back to back,
          with the remaining CRT chunks underneath \\
\bottomrule
\end{tabular}
\end{center}
About \us{470} of the span is PCIe transfer at full speed; what remains is about
\us{95} of kernels and \us{25}--\us{100} of host work at the edges.

\subsection{Point \#1}
Wall clock about \ms{2.3} median natively (first version: \ms{3.09}); GPU span \ms{2.24}.
One call, in \code{hmm} mode:
\begin{center}
\begin{tabular}{rp{0.74\textwidth}}
\toprule
time (\si{\micro\second}) & activity \\
\midrule
0--1440 & input: $a$'s and $b$'s decompose kernels reading the caller's vectors over PCIe
          (about \SI{20}{\giga\byte\per\second}); $a$'s forward NTT on its own stream underneath \\
1440--1684 & \us{244} of kernels: column pass on $b$, \code{k\_ntt\_rows\_mul}, inverse
            column pass (about four memory passes over the residues) \\
1684--2340 & output: the CRT writing the result straight into the output vector over PCIe \\
\bottomrule
\end{tabular}
\end{center}
Host overhead outside the GPU span is about \us{60}. What remains above the floor is the
\us{244} of kernels, plus the gap between HMM access (20--\SI{22}{\giga\byte\per\second}) and
DMA (\SI{24.5}{\giga\byte\per\second}).

%======================================================================
\section{What did not help}
\begin{center}
\begin{tabular}{p{0.42\textwidth}p{0.48\textwidth}}
\toprule
attempt & outcome \\
\midrule
CRT: staging $Y$ in shared memory alone & \us{308} $\to$ \us{232}; the big gain came from
  prime groups \\
More chunks after the copy pool & slower: fixed cost per parallel copy \\
Uneven input chunks & neutral or worse (GPU span \us{625}/\us{647} vs.\ \us{591}--\us{610}) \\
Zero-copy (\code{mapped}) transfers & about 2\,\%, within noise; kept opt-in \\
Spin loop without \code{pause} & no difference \\
Decompose: templated $L$; lazy reduction alone & \us{20.6} $\to$ \us{19.8} $\to$ \us{19.4};
  the limit was table loads \\
Interleaved twiddle table & about \us{1} per NTT kernel; the cost is bytes \\
Generated twiddles at Point \#2 & slower (\code{k\_ntt\_rows\_mul} 31 $\to$ \us{38}); the
  tables fit in L2 there, hence the automatic rule \\
CRT: 2 coefficients per thread; unroll by 8 & fewer blocks / more registers: slower \\
Point \#1 host side: more threads, chunk shapes, later zero-fill & no gain; the later
  zero-fill was worse \\
\bottomrule
\end{tabular}
\end{center}

%======================================================================
\section{Summary at Point \#2}
\begin{center}
\begin{tabular}{lr}
\toprule
stage & median per call \\
\midrule
first GPU run (GPU shared, best setting) & \ms{3.07} \\
same day, old code, idle GPU (best setting) & \ms{2.14} \\
after the copy pool and shallow pipeline & about \us{880} \\
after spinning workers & about \us{760} \\
after helper zero-fill and $1{:}2{:}1$ output chunks & about \us{700} \\
inside FHERMA's challenge image & \us{703} \\
\bottomrule
\end{tabular}
\end{center}
The NTT overlap and the decompose and twiddle changes shortened the GPU span by about
\us{40} more, inside the wall-clock noise.

\appendix
\section{Runtime settings}\label{app:settings}
\begin{center}
\begin{tabular}{lll}
\toprule
variable & default & meaning \\
\midrule
\code{FHERMA\_IN}, \code{FHERMA\_OUT} & \code{hmm} & \code{hmm}, \code{staged}, \code{register}, \code{mapped} \\
\code{FHERMA\_CHUNKS} & 1 (\code{hmm}), 2 & input chunks (equal) \\
\code{FHERMA\_OUT\_CHUNKS} & --- & output chunks (equal) \\
\code{FHERMA\_A/B/C\_CHUNKS} & $C$: \code{1} (\code{hmm}), \code{1,2,1} & chunk weights per operand \\
\code{FHERMA\_THREADS} & $\min(16, \text{cores}/2)$ & host copy threads \\
\code{FHERMA\_SPIN\_MS} & 20 & spin after a call before sleeping \\
\code{FHERMA\_TWIDDLES} & \code{auto} & \code{table}, \code{gen} or \code{auto} (by table size) \\
\code{FHERMA\_PROFILE} & off & per-phase GPU times and host timeline \\
\bottomrule
\end{tabular}
\end{center}

\section{Reproducing the measurements}
\begin{lstlisting}
./build.sh                                   # all binaries (see HANDOVER.md)
./run_tests.sh && ./run_tests.sh build/solution
python3 make_big.py /tmp/p2 32768 868 21     # Point #2
python3 make_big.py /tmp/p1 65536 1747 21    # Point #1
build/solution /tmp/p2 && python3 oracle.py verify /tmp/p2
nsys profile -t cuda --sample=none --cpuctxsw=none -o prof build/solution /tmp/p2
./docker_check.sh 0                          # build and run in fherma/cupqc:0.6.0
\end{lstlisting}

\end{document}
